% Lösungen Einheit 2 von 6 – Nullstellen (zusammenbau v0.1)
\einheitenkopf{Einheit 2 von 6 · Nullstellen}

\erg{12}{a) Faktoren null setzen \quad b) $x_1 = 2$, $x_2 = -5$ \quad c) $x(x^2 - 9) = 0$: $x_1 = 0$, $x_2 = -3$, $x_3 = 3$ \quad d) $x_1 = 1$, $x_2 = 5$ \quad e) $z = 4$ oder $z = 9$: $x_1 = -3$, $x_2 = -2$, $x_3 = 2$, $x_4 = 3$ \quad f) $f(1) = 0$; $f(x) = (x - 1)(x^2 - x - 6)$; $x_1 = 1$, $x_2 = -2$, $x_3 = 3$ \quad g) $e^x > 0$, also $x^2 - 4 = 0$: $x_1 = -2$, $x_2 = 2$ \quad h) $S_y(0 | 0)$; $N_1(-2 | 0)$, $N_2(0 | 0)$, $N_3(3 | 0)$ \quad i) Nullstellen $-3$ und $3$; Breite $= 6$ m \quad j) $S_y(0 | 1{,}8)$; $x^4 - 10x^2 + 9 = 0$, $z = 1$ oder $z = 9$: $N_1(-3 | 0)$, $N_2(-1 | 0)$, $N_3(1 | 0)$, $N_4(3 | 0)$}
\erg{13}{a) $S_y(0 | 5)$; keine Nullstelle: für $x \to \pm\infty$ wächst $f(x)$ über alle Grenzen, der einzige Tiefpunkt liegt über der x-Achse}
\erg{14}{a) Fehler: beim Teilen durch $x$ geht die Lösung null verloren. Richtig: $x(x^2 - 16) = 0$, also $x_1 = 0$, $x_2 = -4$, $x_3 = 4$ \quad b) Ein Produkt ist genau dann null, wenn einer der Faktoren null ist; der erste Faktor ist nur bei $5$ null, der zweite nur bei $-1$. \quad c) $x^2 - 16x - 36 = 0$: $x = 18$ (die Lösung $-2$ entfällt); Weite $18$ m}
